\chapter{Glossário e modelo de estados}\label{ap:glossario}

\section{Termos e siglas}
\begin{quadrosrs}
\caption{Vocabulário do domínio e da especificação}
\begin{tabelasrs}{L{0.24\textwidth}Y}
\cabecalhosrs\textbf{Termo} & \textbf{Definição}\\\midrule
Evento & Atividade com capacidade, instante de início e publicação.\\
Participante & Pessoa identificada pelo provedor externo de identidade.\\
Inscrição ativa & Registro que associa participante e evento e ocupa vaga.\\
Inscrição cancelada & Registro conservado que não ocupa vaga.\\
Precondição & Estado que delimita a execução de um comportamento.\\
Critério de aceite & Resultado observável para avaliar uma obrigação do produto.\\
Idempotência & Repetir uma operação sem acrescentar outro efeito sobre o estado.\\
Rastreabilidade & Relação explícita entre objetivo, origem, requisito e verificação.\\
SRS & Especificação de requisitos de software.\\
UTC & Tempo Universal Coordenado, referência para armazenar instantes.\\
p95 & Percentil 95: limite abaixo do qual ficam pelo menos 95\% das observações.\\
\end{tabelasrs}
\fonte{Elaboração didática deste guia.}
\end{quadrosrs}

\Needspace{0.65\textheight}
\section{Estados da inscrição}
Os estados persistidos são ativa e cancelada. Ausência de inscrição ativa é uma
condição para criar um novo registro, não um terceiro estado armazenado.

\diagramasrs{UML de estados persistidos e transições permitidas}{fig:uml-estados}
  {\begin{tikzpicture}
  \node[uml/estado] (ativa) at (0,0) {Ativa};
  \node[uml/estado] (cancelada) at (7.0,0) {Cancelada};
  \node[circle,fill=black,inner sep=2pt] (inicio) at (-2.6,0) {};
  \draw[uml/fluxo] (inicio) -- (ativa);
  \draw[uml/fluxo] (ativa) -- node[uml/rotulo,above,text width=48mm]
    {cancelar [titular e no prazo]\\ / registrar instante; liberar uma vaga} (cancelada);
  \draw[uml/fluxo] (ativa) edge[loop above,min distance=14mm]
    node[uml/rotulo,above,text width=48mm]
    {repetir criação [mesmo titular/evento]\\ / retornar mesmo ID} (ativa);
  \draw[uml/fluxo] (cancelada) edge[loop below,min distance=14mm]
    node[uml/rotulo,below,text width=52mm]
    {repetir cancelamento [titular e no prazo]\\ / preservar instante e vagas} (cancelada);
\end{tikzpicture}
}

O nó inicial representa o começo do ciclo de vida após uma criação elegível,
sem gatilho ou guarda nessa seta. A elegibilidade da criação está no fluxo de
atividades. As operações pessoais pressupõem sessão válida. As guardas das
transições delimitam prazo e titularidade; os efeitos após a barra indicam
o resultado da transição.
Uma reinscrição cria outro registro e inicia sua própria máquina de estados;
 não reativa a inscrição cancelada. Recusas não mudam o estado existente.


\begin{quadrosrs}
\caption{Condições de transição e efeito nas vagas}
\begin{tabelasrs}{L{0.21\textwidth}YY}
\cabecalhosrs\textbf{Condição inicial} & \textbf{Operação} & \textbf{Resultado}\\\midrule
Sem inscrição ativa & Criar com vaga, publicação e no prazo. & Novo registro ativo; uma vaga consumida.\\
Ativa & Repetir criação para mesmo participante/evento. & Mesmo ID e situação; vagas inalteradas.\\
Ativa & Titular cancela no prazo. & Registro cancelado; uma vaga liberada.\\
Cancelada & Titular repete cancelamento no prazo. & Mesmo instante de cancelamento e vagas.\\
Sem inscrição ativa & Reinscrever após cancelamento, com vaga e no prazo. & Novo ID ativo; registro anterior conservado.\\
Ativa & Cancelar no início ou depois dele. & Operação negada; registro permanece ativo.\\
\end{tabelasrs}
\fonte{Elaboração didática deste guia.}
\end{quadrosrs}
